% Concise closing sections for the 15-page SpaceCHI review manuscript.
% The named manuscript retains the expanded exposition in sections 06--10.
\section{Discussion}
\label{sec:discussion}

The principal finding is a controlled negative benchmark. The runs completed
validly, but Q01 was far behind the physics-residual control, Q01b did not close
the gap, FQK failed its multi-metric screening rule, and no bounded successor
qualified for Gate 6. D046 illustrates why matched controls matter: its 5.80\%
raw improvement over C06 fell to 4.65\% against a classical TWO-RBF model with
the same staged construction, below the predeclared 5\% rule. The inference is
narrow: these feature maps failed under this public-data generator, grouped
development procedure, resource envelope, and comparator set. It does not rule
out all QML or future quantum hardware.

\subsection{Human--AI interaction contribution}

The HCI contribution is an evidence design for decision support, not an
autonomous navigator or a user-tested interface. Before acting on a model
output, an analyst should be able to inspect intended use, data and model
scope, comparators, uncertainty and failure conditions, gate status, and the
action reserved for human authorization. This supports appropriate reliance
on automation \cite{lee2004} and makes capabilities and limitations visible as
recommended in human--AI interaction guidance \cite{amershi2019}.

The ledger also preserves distinctions that a red/green model indicator would
erase. Q02 and Q03 are
\texttt{not\_reached\_under\_frozen\_eligibility}, not failed experiments:
their first-rung tasks ran correctly, but too few configurations met a rule
fixed before outcomes to authorize larger rungs. Valid negatives,
source-ineligible cases, and unopened gates likewise carry different meanings.
Displaying those states can prevent a reviewer from inferring a test or
selection path that never occurred. Whether such a display improves
comprehension, workload, situation awareness, or calibrated reliance remains a
question for a separate study with intended users.

\subsection{From benchmark evidence to mission practice}

Table \ref{tab:evidence-display} turns the governance ledger into requirements
for a future evidence display. They are design implications from this
benchmark, not evaluated usability findings.

The most defensible real-world use of this work is upstream of flight software.
During concept development, an analyst might use a surrogate to prioritize
which candidate correction scenarios receive expensive physical propagation.
The benchmark shows what must accompany such a recommendation: the scenario
and source boundary, comparator performance, uncertainty, numerical-failure
status, and the authorization state. A candidate that passes this offline
stage would still require locked-data evaluation, independent trajectory
analysis, operational constraints, software assurance, and human review before
it could inform a mission decision.

For interface design, this suggests an evidence panel rather than a single
confidence badge. The panel should separate predicted correction cost from
feasibility, identify whether a result is development, calibration, or test
evidence, expose the physical and learned comparators, and explain why an
experiment stopped. It should also preserve a visible fallback to the
validated propagation workflow. In this study, that design would show that
Q01 and Q01b are valid negatives, Q02 and Q03 stopped under eligibility, FQK
failed a multi-metric screening rule, and Gate 6 is unopened. These states lead
to different review actions even though none authorizes deployment.

The societal impact is therefore procedural rather than a claimed improvement
in spacecraft performance. Public agencies, researchers, and reviewers gain a
reproducible way to challenge high-visibility AI claims before scarce compute,
engineering attention, or trust is committed to them. The same pattern can
apply to other high-consequence decision-support domains, but its effect on
human judgment must be tested rather than assumed.

\begin{table}[!t]
\caption{Proposed evidence-display requirements for future human review.}
\label{tab:evidence-display}
\centering
\small
\setlength{\tabcolsep}{4pt}
\renewcommand{\arraystretch}{0.96}
\begin{tabularx}{\textwidth}{@{}p{0.22\textwidth}YY@{}}
\toprule
Field & Information shown & Misinterpretation constrained \\
\midrule
Scope & Sources, force model, uncertainty, and intended use & A development
surrogate is mistaken for mission truth. \\
Evidence stage & Development, calibration, final test, or unopened gate & A
locked or unexecuted evaluation is assumed to have passed. \\
Comparators & Physical residual and matched classical alternatives & A gain
against one weak baseline is attributed to the quantum component. \\
Outcome status & Valid negative, ineligible, source-ineligible, or failed task
& Scientifically different stopping states collapse into one failure label. \\
Authority & Human decision owner and validated propagation fallback &
A prediction is treated as an autonomous command or replaces physical review. \\
\bottomrule
\end{tabularx}
\renewcommand{\arraystretch}{1}
\end{table}

\subsection{Benchmark and safety implications}

A generic classical baseline is inadequate when a low-fidelity physical model
captures target structure directly. Here that residual structure was more
informative than the tested quantum-fidelity features. The matched TWO-RBF
control does not claim to be universally optimal; it removes one credible
classical explanation for D046. The transferable question is whether a
candidate's incremental mechanism remains supported when a classical model
receives the same residual construction and resources.

FQK and CSAFE-RF also show why metrics must be tied to label meaning and human
action. AUROC measures ranking, Brier score probability error, recall the share
of feasible options retained, and false-positive rate the separate risk of
admitting infeasible options. FQK's above-random AUROC cannot offset its failed
Brier and recall criteria. The higher-recall classical result is a future
lesson, not a safety claim: a new protocol must freeze minimum feasible-case
recall, maximum infeasible-case false positives, calibration, and abstention
behavior before fitting.

\section{Limitations and Threats to Validity}
\label{sec:limitations}

The evidence is public-data and development-only: calibration and final-test
rows were not read, so final-test error, operational calibration, mission-loop
behavior, propellant savings, and flight readiness cannot be estimated. Gate 3
supports only eligible public endpoint comparisons; RTC3 is source-ineligible,
and the generated scenarios are not mission-owned truth data. Retaining that
case protects the source boundary but also leaves part of the desired
validation coverage unresolved.

Every quantum run used exact classical statevector simulation at no more than
eight qubits on a laptop. Shot noise, gate and readout errors, device
connectivity, state preparation, transpilation, hardware calibration, and
end-to-end quantum runtime were not measured. Exact simulation removes one
source of sampling variation but cannot establish speedup, hardware efficiency,
or practical quantum advantage.

Finally, frozen preprocessing, feature maps, rungs, controls, and thresholds
make the comparison auditable while restricting inference to the tested
design. The documented literature review is bounded rather than systematic,
so it cannot prove that no related method exists. Alternative models or broader
novelty claims require a new protocol rather than post-outcome modification of
this one.

\section{Future Research Boundary}

Task-specific symmetry and residual structure remain hypotheses, not findings.
A future geometric-QML study must first establish whether any nontrivial
symmetry survives fixed Earth--Moon--Sun ephemerides, thrust constraints, and
reference frames; equivariance alone would not imply advantage over classical
equivariant models. It would then require a mathematical contract, symmetry
audit, baseline reproduction, matched classical comparison, ablations, resource
accounting, and explicit rejection criteria.

D054-A/P018 records bindings for a possible dynamical-Lie-algebra audit, but
closure decision D055-C declined execution. This paper therefore contains no
DOP853 trajectory replay, SVD rank or nullspace result, new symmetry dataset,
variational-circuit training, or geometric-QML performance claim. Any reopening
must occur under a separately authorized protocol able to reject the hypothesis.

\section{Data, Code, and Figure Availability}

Governance records, machine-readable outputs, figure provenance, and source are
available in \href{https://github.com/taechasith/QMLforArtemisIV}{the project
repository} and \href{https://github.com/taechasith/QMLforArtemisIV/releases/tag/v0.3.0}{release
v0.3.0}. Locked rows, mission-owned evidence, trained weights, and unexecuted
future-protocol outputs are not included.

\section{Conclusion}

Under the frozen development benchmark, Q01, Q01b, FQK, and all 16 bounded
successors failed their advancement criteria; none is eligible for Gate 6. The
main contribution is a reproducible negative result and a human-centered claim
structure showing what was tested, why it did not advance, and what remains
subject to human authorization. Its societal value is accountable evaluation
of AI evidence in high-consequence settings, not a deployed mission capability.
The study establishes neither quantum advantage nor operational benefit.
